% BANKS-SOURCE: pdf=183 print=173 kind=section id=9.9
\subsection{Renormalization-group equations in QED}

In this section the physical QED theory is four-dimensional, $D=4$; the
dimension used for dimensional regularization is denoted by
$d_{\mathrm{reg}}$.

% BANKS-SOURCE: pdf=183 print=173 kind=source-page
% BANKS-SOURCE: pdf=183 print=173 kind=prose
We have argued that, in Landau gauge, we can, order by order in the perturbation
expansion, render the expressions for all Green functions finite by rescaling the
fields and tuning the bare parameters $m_0$ and $e_0$ so that the renormalized
parameters $m=Z_m m_0$ and $e$ are finite. In the process, we have introduced a
new mass scale, $\mu$, the renormalization scale, writing the dimensional bare
coupling $e_0^2$ as
% BANKS-SOURCE: pdf=183 print=173 kind=equation
\[
  e_0^2=\mu^{4-d_{\mathrm{reg}}}e^2Z_3^{-1}.
\]
The Green functions of the bare fields, expressed as functions of $e_0$ and
$m_0$, are independent of $\mu$:
% BANKS-SOURCE: pdf=183 print=173 kind=equation
\[
  \mu\frac{\dd}{\dd\mu}\,\Gamma^{(0)}_{F,A}(e_0,m_0)=0.
\]
$\Gamma^{(0)}_{F,A}$ is the 1PI vertex with $F$ fermion and $A$ photon legs.
It is a function of the momenta of the external legs, defined without the
momentum-conserving delta function. It has mass dimension $4-\frac32F-A$. The
renormalized Green function is
% BANKS-SOURCE: pdf=183 print=173 kind=equation
\[
  \Gamma_{F,A}=Z_2^{F/2}Z_3^{A/2}\Gamma^{(0)}_{F,A}.
\]
$\mu$-independence of the bare vertices implies a relation for the renormalized
ones. It should be expressed in terms of the renormalized parameters, and we
can do this using the chain rule:
% BANKS-SOURCE: pdf=183 print=173 kind=equation
\[
  \left(\mu\partial_\mu+\beta\partial_\alpha
  +\gamma_A\,\kappa\partial_\kappa+m\gamma_m\partial_m\right)\Gamma_{F,A}
  =-(F\gamma_F+A\gamma_A)\Gamma_{F,A},
\]
where
% BANKS-SOURCE: pdf=183 print=173 kind=equation
\[
  \beta=\mu\frac{\dd}{\dd\mu}\alpha,
  \qquad
  \gamma_m=\mu\frac{\dd}{\dd\mu}\ln Z_m,
  \qquad
  \gamma_F=\frac12\mu\frac{\dd}{\dd\mu}\ln Z_2,
  \qquad
  \gamma_A=\frac12\mu\frac{\dd}{\dd\mu}\ln Z_3.
\]
We have used the symbol $\dd/\dd\mu$ to represent derivatives with $\alpha_0$
and $m_0$ held fixed, while $\partial_\mu$ refers to derivatives with the
renormalized parameters fixed. $\alpha_0=e_0^2/(4\pi)$ is the bare
fine-structure constant. The renormalized fine-structure constant is
% BANKS-SOURCE: pdf=183 print=173 kind=equation
\[
  \alpha=\mu^{d_{\mathrm{reg}}-4}\alpha_0Z_3.
\]
Thus
% BANKS-SOURCE: pdf=183 print=173 kind=equation
\[
  \beta=(d_{\mathrm{reg}}-4)\alpha+2\gamma_A.
\]
We have also used the result that the longitudinal part of the photon
propagator is unchanged by loop corrections. Thus, the renormalization of the
gauge parameter $\kappa$

% BANKS-SOURCE: pdf=184 print=174 kind=source-page
comes only from the photon wave-function renormalization $\kappa_0=Z_3\kappa$.
This equation for the renormalized 1PI vertices is called the
renormalization-group equation. It holds for all values of momenta, and
expresses the fact that the renormalized theory has only one independent
dimensionless parameter and we have artificially introduced another one
through the mass scale $\mu$. As we have emphasized, the reason why we must do
this is because the theory is not scale-invariant even in the limit $m_0\to0$.
If we had insisted on defining the renormalization scale in terms of the
electron mass parameter, we would have introduced spurious infrared
divergences into the $m_0=0$ theory.

The fact that the RG equation is true for all momenta shows us that the RG
functions $\beta$ and $\gamma_i$ are all finite in the limit
$d_{\mathrm{reg}}\to4$. Read as
a set of equations for these functions in terms of the finite proper vertices,
we have an over-determined set of linear equations for $\beta$ and $\gamma_i$.
The fact that they have a solution is the content of the differential equation
for the vertices. In perturbation theory, it is easy to see that, to $k$th
order in the loop expansion, we can expect the functions $Z_i$ to have poles up
to order $k$ at $d_{\mathrm{reg}}=4$. Expressed in terms of $\alpha_0$, the
$\mu$ dependence of the $k$th order is just
$\mu^{k(d_{\mathrm{reg}}-4)}$. The scaling derivative in the definition of the
RG functions brings down a factor of $d_{\mathrm{reg}}-4$, which cancels
out the first-order pole. All the others must cancel out automatically, when
the finite functions $\beta$ and $\gamma_i$ are expressed in terms of $\alpha$,
$\kappa$, and $m$! Thus, to calculate the $k$th-order term in any of these
functions, we need only find the residue of the first-order pole, in the $k$th
order in the loop expansion.

\BanksImplicitHook{I-CH09-017}

Finally we note that, in the MS prescription, the coefficients of these poles
are independent of $m_0$. This follows from the fact that all of the $Z_i$ are
dimensionless and therefore $(\partial/\partial m_0)\ln Z_i$ is given by
convergent Feynman integrals in $D=4$, plus the fact that we have introduced
$\mu$ rather than $m_0$ or $m$ to provide the dimensions of $e_0$. It follows
that all the RG functions depend only on the renormalized coupling $\alpha$ and
the gauge parameter $\kappa$. In fact, since the transverse part of the photon
propagator is $\kappa$-independent, $\gamma_A$ and $\beta$ are independent of
$\kappa$. In a more general field theory, defined in perturbation theory around
a Gaussian fixed point, the RG functions in the MS scheme depend on the
marginal parameters, but not on the relevant ones.

If we rescale all the momenta $p_i\to e^t p_i$, then the equation of
dimensional analysis is
% BANKS-SOURCE: pdf=184 print=174 kind=equation
\[
  \left(\partial_t+\mu\partial_\mu+m\partial_m\right)\Gamma_{F,A}
  =\left(4-\frac32F-A\right)\Gamma_{F,A}.
\]
Using the RG equation, we can rewrite this as
% BANKS-SOURCE: pdf=184 print=174 kind=equation
\[
  \left(\partial_t+\beta\partial_\alpha
  +m(1+\gamma_m)\partial_m\right)\Gamma_{F,A}
  =-\left(\frac32F+\gamma_F+(1+\gamma_A)A\right)\Gamma_{F,A}.
\]
In words, this equation says that we can compensate for a change in momentum
scale by changing the coupling and the mass, and rescaling the fields. Thus the
solution of this equation is
% BANKS-SOURCE: pdf=184 print=174 kind=equation
\[
  \Gamma_{F,A}(e^t p_i,\alpha,m)
  =\Gamma_{F,A}(p_i,\alpha(t),m(t))e^{L(t)}.
\]

% BANKS-SOURCE: pdf=185 print=175 kind=source-page
On differentiating this expression with respect to $t$ and insisting that it
give the previous equation we find that
% BANKS-SOURCE: pdf=185 print=175 kind=equation
\[
  \dot\alpha(t)=\beta,
  \qquad
  \dot m(t)=(1+\gamma_m)m,
  \qquad
  \dot L(t)=4-F\Theta_F-A\Theta_A,
\]
with $\Theta_F=\frac32+\gamma_F$ and $\Theta_A=1+\gamma_A$.

Thus, we learn that a rescaling of all momenta is equivalent to a flow of the
coupling constants satisfying the above equations, combined with a
scale-dependent rescaling of the Green functions, which can be interpreted as
an anomalous dimension for the fields. Similarly, the flow equation for the
mass can be interpreted as an anomalous dimension for the renormalized mass
parameter. The anomalous dimensions $\Theta_{F,A,m}=E_{F,A,m}+\gamma_{F,A,m}$,
where $E_{F,A,m}=(3/2,1,1)$ is the engineering dimension,\footnote[13]{Often
the term anomalous dimension is reserved for the shift $\gamma_{F,A,m}$ rather
than the dimension itself.} are functions of the scale-dependent
fine-structure constant. If, in the asymptotic region $t\to\pm\infty$,
$\alpha(t)\to\alpha^*$, a fixed point, then we conclude that the theory is
asymptotically scale-invariant, with dimensions given by the anomalous
dimensions evaluated at the fixed point.

Are there fixed points in QED? We can investigate this only in the vicinity of
the free theory, because we are doing perturbation theory. The perturbative
value of the $\beta$ function in spinor QED is
% BANKS-SOURCE: pdf=185 print=175 kind=equation
\[
  \beta=(d_{\mathrm{reg}}-4)\alpha+\frac{2}{3\pi}\alpha^2.
\]
Adding scalars and more spinors changes the coefficient of the second term,
without changing its sign. We note that for real values of $d_{\mathrm{reg}}$
less than but
close to $4$ we find a zero of $\beta$ at a non-zero but small value of the
coupling, where we can trust perturbation theory. This will occur in any theory
with couplings that are dimensionless in four dimensions but have positive
coefficient for the one-loop correction to the RG function. This observation
is the basis of one of the methods of calculating critical exponents for
second-order phase transitions in Euclidean dimensions $d_{\mathrm{reg}}=2,3$.
One computes the fixed points and anomalous dimensions in a power series in
$4-d_{\mathrm{reg}}$, and extrapolates to the values
of the dimension relevant for real condensed-matter systems. This turns out to
give quite good agreement with experiment, though other methods based on
field-theory calculations in integer dimensions do somewhat better. A more
detailed account can be found in the books of Peskin and Schroeder and
Zinn--Justin \cite{banks-ref-33,banks-ref-111} and the references to the
original literature found therein.

As high-energy physicists, we are interested in $D=4$. There, the RG equation
has the solution
% BANKS-SOURCE: pdf=185 print=175 kind=equation
\[
  \alpha(t)=\frac{\alpha(0)}{1-[2/(3\pi)]\alpha(0)t}.
\]
Note that, for $t\to-\infty$, this solution remains in the perturbative regime
if $\alpha(0)$ is small. On the other hand, if $t\to\infty$ the coupling
becomes strong, and in fact appears

% BANKS-SOURCE: pdf=186 print=176 kind=source-page
to reach infinity at the Landau pole, $t_L^{-1}=[2/(3\pi)]\alpha(0)$. Our
basic definition of field theory tells us that a mathematically consistent
field theory will approach its fixed point in the deep UV. Our perturbative
philosophy was based on the assumption that this fixed point was Gaussian. The
Landau pole tells us that this assumption was wrong. The QED coupling is
marginally irrelevant, which means that the only mathematically consistent
value for the renormalized coupling is zero. If we insist on a finite value for
this coupling at some momentum scale $t=0$, then the theory becomes singular at
a finite scale, that of the Landau pole.

Since $\alpha(0)\sim1/137$ in the real world, our attitude to this as physicists
is somewhat different than that of the mathematical field theorist. If $t=0$
corresponds to the atomic scale $\sim10\,\mathrm{eV}$ at which the
fine-structure constant is measured, then the scale of the Landau pole is
$\sim10^{331}\,\mathrm{GeV}$, which is much larger than the Planck scale
$10^{19}\,\mathrm{GeV}$ defined by quantum gravity. We certainly expect our
notions of quantum field theory to break down long before we reach the scale of
the Landau pole. Thus, QED is a perfectly good effective field theory at all
scales at which we expect the whole notion of quantum field theory to work. In
contrast with a truly irrelevant operator, like the four-fermion coupling of
Fermi's theory of weak interactions, a small, marginally irrelevant coupling
does not hint at new physics around the corner. The renormalized perturbation
series of QED could in fact be the whole story up to the Planck scale. We know
that this is not true, because, above the QCD and weak-interaction scales,
physics becomes more complicated. QED is certainly incorporated into the full
standard model. The point is that there is nothing in pure QED that could have
led us to such a conclusion.

% BANKS-SOURCE: pdf=186 print=176 kind=subsection id=9.9.1
\subsubsection{The static potential and the definition of $\alpha$}

% CONVENTION-SCOPE: fixed-D=4 reason=physical QED Wilson line

% BANKS-SOURCE: pdf=186 print=176 kind=prose
The flip side of the marginal irrelevance of QED is that the coupling becomes
arbitrarily weak in the IR. Indeed, an alternative name for marginally
irrelevant interactions is IR-free. One must, however, be careful about the
precise meaning of this. In the IR limit, the mass parameter also becomes
large. It turns out that this makes the MS definition of the coupling very
different from a physical definition in terms of real scattering amplitudes.
One useful physical definition is in terms of the long-distance potential felt
by a pair of oppositely charged heavy particles. If a particle is sufficiently
heavy we can describe its interaction with the electromagnetic field in terms
of its classical trajectory through space-time, $x^\mu(\tau)$. The current of
the particle is
% CONVENTION-SCOPE: fixed-D=4 reason=physical QED worldline current
% BANKS-SOURCE: pdf=186 print=176 kind=equation
\[
  J^\mu(x)=q\int\dd\tau\,\frac{\dd x^\mu}{\dd\tau}
  \delta^4(x-x(\tau)),
\]
where $q$ is its charge. The interaction with the electromagnetic field is
obtained by inserting the factor
% CONVENTION-SCOPE: fixed-D=4 reason=physical QED Wilson-line interaction
% BANKS-SOURCE: pdf=186 print=176 kind=equation
\[
  \ee^{\ii\int\dd^4x\,A_\mu(x)J^\mu(x)}
  =\ee^{\ii q\int A_\mu\,\dd x^\mu},
\]
where the second integral is a line integral along the particle path. This
expression is the Wilson line. A similar formula holds in non-abelian gauge
theory, where $A_\mu$ is a matrix and the exponential is replaced by a
path-ordered exponential along the particle path.

% BANKS-SOURCE: pdf=187 print=177 kind=source-page
It is convenient to discuss a particle--anti-particle pair as a closed Wilson
line, or Wilson loop. This corresponds to pair production, propagation through
space-time, and annihilation at some later time. The pair-production and
annihilation events are described in an idealized manner and do not correspond
to a realistic experiment. However, if we take a rectangular loop in Euclidean
space with one side $T$ much greater than the other, $T\gg R$, and Wick rotate
to Lorentzian signature, then to leading order in $T/R$ the answer will be
dominated by the long period of particle--anti-particle propagation, and the
unrealistic pair creation events will be a sub-leading effect.

The exact answer in Euclidean QED is given in terms of the generating function
of connected Green functions, by $\ee^{-W[J]}$. However, in the large-$R$ limit,
the contribution of the two-point function dominates (cluster decomposition
works even in this theory with massless photons). Furthermore, in the same
limit, the contribution of the two-point function is dominated by its behavior
at $p=0$. If we assume that $e_0^2/[1+\Pi(p^2)]$ goes to a finite $p=0$ limit,
then the Wilson loop is given by
% BANKS-SOURCE: pdf=187 print=177 kind=equation
\[
  \ee^{-T\,\frac{e_0^2}{1+\Pi(0)}\,\frac{1}{4\pi R}}.
\]
This is the Coulomb potential with physical renormalized charge
% BANKS-SOURCE: pdf=187 print=177 kind=equation
\[
  e_{\mathrm{Coul}}^2=\frac{e_0^2}{1+\Pi(0)}.
\]
We have already noted that we could have chosen to parametrize the renormalized
theory by this quantity. However, we also commented that this introduced a
logarithmic dependence on $m_0$ (and thus on $m$) and we preferred to introduce
the arbitrary momentum scale $\mu$ to eliminate this. For finite $m$, this
choice of parametrization cannot change the fact that $e_{\mathrm{Coul}}$ is a
finite quantity. Thus, the physical renormalized charge defined in terms of the
Coulomb potential is not equal to the IR limit of the MS coupling
$\alpha(-\infty)=0$. Physically, this corresponds to the fact that the
renormalization of the electric charge due to loops of electrons goes to zero
below the threshold for electron--positron pair production. Mathematically it
corresponds to the fact that one must take $m(t)$ to infinity as $\alpha(t)$
goes to zero in computing $e_{\mathrm{Coul}}$ in the MS scheme.\footnote[14]{IR
freedom is useful in the massive theory when dealing with the problem of
soft-photon IR divergences. It justifies the perturbative resummation
procedures one uses to extract finite answers for inclusive cross sections.}

In the $m=0$ theory, on the other hand, the MS RG equation can be used to show
that the Fourier transform of the ``Coulomb'' potential is modified to
$1/[q^2\ln(q^2)]$, so that the potential falls off more rapidly than $1/R$,
corresponding to vanishing charge in the IR.
